\section{Previous work and positioning}
\label{sec:prevwork}

\Cref{tab:prior-art} lists, element by element, which parts of the
procedure are established prior results reused unchanged and which parts
are this report's own methodological or policy contribution. Nothing in
\cref{sec:model,sec:freq} is claimed as new physics.

\begin{table*}[htbp]
\centering
\caption{Established prior results integrated by this report vs. its own methodological and policy contribution.}
\label{tab:prior-art}
\begin{tabular}{p{5.1cm}p{2.6cm}p{6.7cm}}
\toprule
element & class & this report's treatment\\
\midrule
Sealed-box alignment, Thiele/Small closed-box theory & established prior result & reused unchanged as the sealed-box ODE and alignment line\\
Displacement-limited output ceiling ($V_d=S_dX_{\max}$) & established prior result & used as an exact descriptor, not a universal loudspeaker optimum\\
Thermal (dissipation-limited) output ceiling ($\eta_0P_{\max}$) & established prior result & used as an exact descriptor; never summed with the displacement ceiling\\
$X_{\max}$ as an IEC/Klippel excursion anchor & established prior result & used only as the clipping boundary; not converted into absolute THD\\
Room pressure-zone compliance below the first mode & established prior result & extended with an explicit first-order leakage-corner term; used only as a first-pass sizing/limiting model, not a claim about the complete room\\
Doppler (FM) sideband index & established prior result & reported as a separate, non-summed nonlinear-risk indicator\\
Finite-band DSP room correction & established prior result & adopted as an explicit scope input rather than an unconstrained assumption\\
Role-specific, non-compensatory ranking policy (mono sub manifold vs. independent stereo upper-bass channels) & this report's contribution & declared engineering policy; separate Pareto fronts per role, no cross-role or cross-channel summing credit\\
Explicit finite-amplifier feasibility gates (voltage, current, continuous and burst power) & this report's contribution & integrated into every role evaluation and the worked example\\
Numerical sealed-ODE transient (burst) displacement factor & this report's contribution & replaces a constant free-mass burst factor near $F_c$\\
Aggregate, hash-verified private-database screening evidence & this report's contribution & counts and hexbin/histogram population evidence only, no row-level redistribution\\
\bottomrule
\end{tabular}
\end{table*}


\textbf{Alignment and output ceilings.} Closed-box small-signal analysis
and synthesis \cite{small1972} and direct-radiator system analysis
\cite{small1972direct} supply the sealed-box model, the alignment relation,
the reference efficiency $\eta_0$, and the displacement- and
dissipation-limited output ceilings used throughout. Fielder and Benjamin
\cite{fielder1988} established the corresponding program-material and
room-aware view of subwoofer displacement requirements, which is the
closest prior work to the sizing part of the present procedure.

\textbf{Ratings and their comparability.} Continuous power ratings follow
different standards \cite{aes2}, and published $\Xmax$ values may be
geometric, distortion-defined, or vendor-specific
\cite{iec62458,klippel2003xmax,klippel2006}. Both introduce enough spread
between manufacturers to change a ranking, which is why the procedure
carries these as declared uncertainty inputs
(\cref{sec:scope}) and reports rank robustness against them
(\cref{sec:procedure}).

\textbf{$EBP$ as a screening heuristic.} $EBP\lesssim\SI{50}{\hertz}$
suggesting sealed and $\gtrsim\SI{100}{\hertz}$ suggesting vented
\cite{dickason2006} is a widely used enclosure-\emph{type} screen. It is a
motor-to-mass ratio and, through the exact corner rate $F_s/\Qts$, fixes
where a chosen $\Qtc$ places the box corner; it does not state achievable
output or nonlinear headroom. Here it enters only through the exact
$F_s/\Qts$ used in \cref{eq:alignment}, never as a standalone ranking key.
The approximation $F_s/\Qts\approx EBP$ is not used numerically: it differs
by $(1+\Qes/\Qms)$, which reaches several per cent to over ten per cent for
typical bass drivers.

\textbf{Why role specialization is worth testing.} A limited overhung-motor
scaling follows from the same standard identities. Let $u=h_c/h_g>1$ be
the winding-height/effective-gap-height ratio, let total conductor length
be $\ell$ and area $A_w$, and assume uniform gap flux $B$, conductor
conductivity $\sigma_c$, density $\rho_c$, and coil mass
$m_c=\rho_c\ell A_w$. With only the fraction $1/u$ of the winding in the
effective gap,
\begin{equation}
	\begin{gathered}
		\Bl\simeq\frac{B\ell}{u},\qquad
		\Re=\frac{\ell}{\sigma_cA_w},\qquad
		m_c=\rho_c\ell A_w,\\
		\boxed{\ \frac{F_s}{\Qes}u^2
		=\frac{B^2\sigma_c}{2\pi\rho_c}\frac{m_c}{\Mms}\ } .
	\end{gathered}
	\label{eq:motor-scaling}
\end{equation}
Thus, at fixed flux, material and coil-mass allocation, obtaining more
geometric overhang tends to trade against electrical damping per moving
mass. This is motivation to \emph{compare} a high-stroke low-bass role
with a lower-excursion upper-bass role. It is not a motor bound: $B$,
effective gap, $u$, coil-mass fraction, fringing and topology are normally
absent from datasheets, and expensive high-flux motors can move the
tradeoff. The relation is therefore never a feasibility gate or ranking
key. The architecture comparison of \cref{sec:procedure-architecture}
tests the practical consequence directly and retains any full-band
candidate that passes the same system gates.

\textbf{Large-signal behaviour and nonlinear risk.} Motor and suspension
nonlinearity, and the parameters used to characterize it, are surveyed by
Klippel \cite{klippel2006}; $\Xmax$ definitions and their measurement are
treated in \cite{klippel2003xmax}. Doppler (FM) distortion in
direct-radiator loudspeakers was described by Klipsch \cite{klipsch1969},
and its magnitude and audibility were subsequently disputed by Allison and
Villchur \cite{allison1982}; the present report therefore reports the
small-index first-sideband ratio and cites that debate rather than
converting it into a perceptual penalty. Thermal compression and power
handling follow Button \cite{button1992}. Perceptual thresholds for low-frequency
loudspeaker distortion depend on order, spectrum, level, and program
material \cite{fielder2019}, so no absolute audibility claim is made here.

\textbf{Ports.} Port compression, turbulence, and nonlinearity are
characterized in \cite{salvatti2002,vanderkooy1998}. These references
support the scoped discussion in \cref{app:vented}, which motivates the
choice of sealed loading for this application without claiming general
superiority.

\textbf{Room and correction.} The distinction between direct/early-field
and modal-region behaviour, and the interaction of loudspeakers with small
rooms, follows Toole \cite{toole2006}. Optimization over a defined seating
area using multiple subwoofers \cite{welti2006} is the explicit contrast to
the compact-area objective adopted here. Modal equalization
\cite{makivirta2003} and the review of room-response equalization
\cite{cecchi2018} bound what finite-band correction can achieve, and
motivate treating DSP correction as band-limited and regularized rather
than as exact inversion.

\textbf{Spatial attributes.} Scene-based spatial-quality terminology
\cite{rumsey2002} separates image and scene attributes from environmental
envelopment; small-room reflection effects are quantified in
\cite{bech1998}; envelopment's dependence on late lateral energy
\cite{bradley1995} and listener preference across reproduction methods
\cite{francombe2017} are the counterpoint that keeps
\cref{sec:intro-fidelity}'s claim restricted to scene fidelity.

What is missing from this body of work, and what this report supplies, is
not a new physical relation but a complete, reproducible, datasheet-only
decision path from a candidate population to a ranked, gate-checked
shortlist for one declared architecture, including the finite amplifier,
the finite enclosure, the room, and the transient case.
